\documentclass[../../main.tex]{subfiles}
\begin{document}
\section{分布与对合性}\label{section:ch19-distributions-involutivity}

设 $M$ 是一个光滑流形。$M$ 上一个秩为 $k$ 的\emph{分布}是指 $TM$ 的一个秩-$k$ 子丛。若它是一个光滑子丛，则称为\emph{光滑分布}。分布有时也称为\emph{切分布}（尤其是在可能与分析中把广义函数称为分布这一用法相混淆时）、$k$-平面场或切子丛。

通常，一个秩-$k$ 分布可这样描述：对每个 $p\in M$ 指定一个 $k$ 维线性子空间 $D_p\subseteq T_pM$，并令 $D=\bigcup_{p\in M}D_p$。于是由子丛的局部标架判别法（引理\cref{lemma:ch10-子丛局部标架判别}）可知，$D$ 是光滑分布当且仅当 $M$ 的每一点都有一个邻域 $U$，在其上存在光滑向量场 $X_1,\dots,X_k\colon U\to TM$，使得对每个 $q\in U$，$X_1|_q,\dots,X_k|_q$ 构成 $D_q$ 的一组基。此时称 $D$ 是（局部地）由向量场 $X_1,\dots,X_k$ 张成的分布。

\subsection{积分流形与对合性}

设 $D\subseteq TM$ 是一个光滑分布。若对每个 $p\in N$ 都有 $T_pN=D_p$，则称非空浸入子流形 $N\subseteq M$ 为 $D$ 的一个\emph{积分流形}。本章所要讨论的主要问题就是积分流形的存在性。

在进入一般理论之前，先描述一些应当牢记的分布与积分流形的例子。

\begin{example}[分布与积分流形]\label{example:ch19-分布与积分流形}
(a) 若 $V$ 是流形 $M$ 上一个处处非零的光滑向量场，则 $V$ 在 $M$ 上张成一个光滑的秩-$1$ 分布（见例\cref{example:ch10-子丛例子}）。$V$ 的任意积分曲线的像都是 $D$ 的积分流形。
(b) 在 $\mathbb{R}^n$ 中，向量场 $\partial/\partial x^1,\dots,\partial/\partial x^k$ 张成一个秩为 $k$ 的光滑分布。平行于 $\mathbb{R}^k$ 的 $k$ 维仿射子空间都是积分流形。
(c) 设 $R$ 是 $\mathbb{R}^n\setminus\{0\}$ 上由单位径向向量场 $x^i\,\partial/\partial x^i$ 张成的分布，$R^\perp$ 是其正交补丛（见引理\cref{lemma:ch10-正交补丛}）。则 $R^\perp$ 是 $\mathbb{R}^n\setminus\{0\}$ 上一个光滑的秩-$(n-1)$ 分布。对每个 $x\in\mathbb{R}^n\setminus\{0\}$，以 $0$ 为中心、半径为 $|x|$ 的球面都是 $R^\perp$ 的积分流形。
(d) 设 $D$ 是 $\mathbb{R}^3$ 上由下列向量场张成的光滑分布：
\begin{align*}
X=\frac{\partial}{\partial x}+y\frac{\partial}{\partial z},\qquad Y=\frac{\partial}{\partial y}.
\end{align*}
（见图\cref{figure:ch19-no-integral}。）事实上 $D$ 没有积分流形。为了理解其原因，设 $N$ 是一条过原点的积分流形。由于 $X$ 与 $Y$ 都与 $N$ 相切，从 $N$ 出发的 $X$ 或 $Y$ 的任意积分曲线都必须停留在 $N$ 中，至少在一小段时间内如此（习题\cref{exercise:ch09-02}）。于是 $N$ 含有 $x$ 轴的一个开子集（$x$ 轴是 $X$ 的一条积分曲线）。对每个充分小的 $x$，它还含有过 $(x,0,0)$、平行于 $y$ 轴的直线的一个开子集（该直线是 $Y$ 的一条积分曲线）。因此 $N$ 含有 $(x,y)$ 平面的一个开子集。然而，在 $x$ 轴之外任意一点 $p$ 处，$(x,y)$ 平面的切平面都不等于 $D_p$。因而不存在这样的积分流形。
\end{example}

\begin{figure}[ht]
\centering
\includegraphics[width=0.8\textwidth]{fig-19-1.png}
\caption{无积分流形的光滑分布}\label{figure:ch19-no-integral}
\end{figure}

最后一个例子表明，一般而言积分流形可能不存在。设 $D$ 是 $M$ 上的一个光滑分布。若给定 $D$ 的任意一对光滑局部截面（即定义在 $M$ 的某个开子集上、且对每点 $p$ 满足 $X_p,Y_p\in D_p$ 的光滑向量场 $X,Y$），它们的 Lie 括号仍是 $D$ 的局部截面，则称 $D$ 是\emph{对合的}。下一个命题表明，对合性条件可用 Lie 代数简洁地表述。

\begin{proposition}\label{proposition:ch19-对合即Lie子代数}
设 $D\subseteq TM$ 是一个光滑分布，$\Gamma(D)\subseteq\mathfrak X(M)$ 表示 $D$ 的光滑整体截面空间。则 $D$ 对合当且仅当 $\Gamma(D)$ 是 $\mathfrak X(M)$ 的一个 Lie 子代数。
\end{proposition}

\begin{proof}
若 $D$ 对合，则由定义可知 $\Gamma(D)$ 对 Lie 括号封闭。又因它是 $\mathfrak X(M)$ 的一个线性子空间，故它是一个 Lie 子代数。

反之，设 $\Gamma(D)$ 是 $\mathfrak X(M)$ 的一个 Lie 子代数，$X,Y$ 是 $D$ 在某个开子集 $U\subseteq M$ 上的光滑局部截面。给定 $p\in M$，取 $\psi\in C^\infty(M)$ 为一个鼓包函数，它在 $p$ 的一个邻域上恒等于 $1$，且支集含于 $U$ 中。于是 $\psi X$ 与 $\psi Y$ 都是 $D$ 的光滑整体截面，故由假设它们的 Lie 括号也是 $D$ 的截面。这个 Lie 括号为 $[\psi X,\psi Y]=\psi^2[X,Y]+(X\psi)Y-(Y\psi)X$，它在 $p$ 的一个邻域上等于 $[X,Y]$。因而对每个 $p\in U$ 都有 $[X,Y]_p\in D_p$，故 $D$ 对合。
\end{proof}

若 $M$ 的每一点都含于 $D$ 的某个积分流形中，则称 $M$ 上的光滑分布 $D$ 是\emph{可积的}。

\begin{proposition}\label{proposition:ch19-可积蕴含对合}
每个可积分布都是对合的。
\end{proposition}

\begin{proof}
设 $D\subseteq TM$ 是一个可积分布。设 $X,Y$ 是定义在某个开子集 $U\subseteq M$ 上的 $D$ 的光滑局部截面。令 $p$ 为 $U$ 中任意一点，$N$ 为含 $p$ 的一个 $D$ 的积分流形。$X$ 与 $Y$ 是 $D$ 的截面这一事实意味着 $X$ 与 $Y$ 都与 $N$ 相切。由推论\cref{corollary:ch08-32}，$[X,Y]$ 也与 $N$ 相切，从而 $[X,Y]_p\in D_p$。由于这对每个 $p\in U$ 都成立，故 $D$ 对合。
\end{proof}

例如注意到，例\cref{example:ch19-分布与积分流形}(d) 中的分布 $D$ 不是对合的，因为 $[X,Y]=\partial/\partial z$ 不是 $D$ 的截面。

下一个引理表明，对合性条件不必对每一对光滑向量场都加以检验，只需检验在每一点的一个邻域内的一组光滑局部标架即可。

\begin{lemma}[对合性的局部标架判别法]\label{lemma:ch19-局部标架对合判别}
设 $D\subseteq TM$ 是一个分布。若在 $M$ 的每一点的一个邻域内都存在 $D$ 的一组光滑局部标架 $(V_1,\dots,V_k)$，使得对每个 $i,j=1,\dots,k$ 都有 $[V_i,V_j]$ 是 $D$ 的截面，则 $D$ 对合。
\end{lemma}

\begin{proof}
设假设成立，$X,Y$ 是 $D$ 在某个开子集 $U\subseteq M$ 上的光滑局部截面。给定 $p\in U$，在 $p$ 的一个满足假设的邻域内取一组光滑局部标架 $(V_1,\dots,V_k)$，并在该邻域内写成 $X=X^iV_i$、$Y=Y^jV_j$。于是利用\eqref{align:ch08-08}，
\begin{align*}
[X,Y]=[X^iV_i,Y^jV_j]=X^iY^j[V_i,V_j]+X^i(V_iY^j)V_j-Y^j(V_jX^i)V_i.
\end{align*}
由假设，最后这一式是 $D$ 的截面。
\end{proof}

\subsection{对合性与微分形式}

微分形式为描述分布与对合性提供了另一种方式。

\begin{lemma}[光滑分布的 $1$-形式判别法]\label{lemma:ch19-1形式判别}
设 $M$ 是一个光滑 $n$ 维流形，$D\subseteq TM$ 是一个秩为 $k$ 的分布。则 $D$ 光滑当且仅当每一点 $p\in M$ 都有一个邻域 $U$，在其上存在光滑 $1$-形式 $\omega^1,\dots,\omega^{n-k}$，使得对每个 $q\in U$，
\begin{align}
D_q=\operatorname{Ker}\omega^1\big|_q\cap\cdots\cap\operatorname{Ker}\omega^{n-k}\big|_q.
\end{align}
\end{lemma}

若 $D$ 是光滑 $n$ 维流形 $M$ 上的一个秩-$k$ 分布，在开子集 $U\subseteq M$ 上定义、并对每个 $q\in U$ 满足上式的任意 $n-k$ 个线性无关的 $1$-形式 $\omega^1,\dots,\omega^{n-k}$，称为 $D$ 的一组\emph{局部定义形式}。更一般地，若 $0\leqslant p\leqslant n$，当 $X_1,\dots,X_p$ 为 $D$ 的局部截面时恒有 $\omega(X_1,\dots,X_p)=0$，则称 $p$-形式 $\omega\in\Omega^p(M)$ \emph{零化} $D$（在 $p=0$ 的情形，只有零函数零化 $D$）。

\begin{lemma}\label{lemma:ch19-零化形式结构}
设 $M$ 是一个光滑 $n$ 维流形，$D$ 是 $M$ 上一个光滑的秩-$k$ 分布。设 $\omega^1,\dots,\omega^{n-k}$ 是 $D$ 在开子集 $U\subseteq M$ 上的一组光滑局部定义形式。定义在 $U$ 上的一个光滑 $p$-形式 $\eta$ 零化 $D$ 当且仅当它可写成
\begin{align}
\eta=\sum_{i=1}^{n-k}\omega^i\wedge\beta^i
\end{align}
的形式，其中 $\beta^1,\dots,\beta^{n-k}$ 是 $U$ 上的某些光滑 $(p-1)$-形式。
\end{lemma}

\begin{remark}
在 $p=1$ 的情形，$\beta^i$ 是光滑函数，此时把与光滑函数的外积理解为普通乘法。因而在该情形下，本引理只是说 $\eta$ 是 $\omega^i$ 的以光滑函数为系数的线性组合。
\end{remark}

\begin{proof}
容易验证，在每一点的一个邻域内满足上式的任意形式 $\eta$ 都零化 $D$。反之，设 $\eta$ 在 $U$ 上零化 $D$。在每一点的一个邻域内，可把 $(n-k)$ 元组 $(\omega^1,\dots,\omega^{n-k})$ 补成 $M$ 的一组光滑局部余标架 $(\omega^1,\dots,\omega^n)$（命题\cref{proposition:ch10-扩充为局部标架}）。若 $(E_1,\dots,E_n)$ 是对偶标架，则 $D$ 局部地由 $E_{n-k+1},\dots,E_n$ 张成。用这组余标架，任意 $\eta\in\Omega^p(M)$ 都可局部唯一地写成
\begin{align*}
\eta=\sum_I \eta_I\,\omega^{i_1}\wedge\cdots\wedge\omega^{i_p},
\end{align*}
其中系数由 $\eta_I=\eta(E_{i_1},\dots,E_{i_p})$ 确定。于是 $\eta$ 零化 $D$ 当且仅当凡是含有某个 $\omega^{n-k+1},\dots,\omega^n$ 因子的系数 $\eta_I$ 全为零，这等价于 $\eta$ 具有上述形式。
\end{proof}

下面的命题用外导数给出对合性的另一种刻画。

\begin{proposition}\label{proposition:ch19-对合外导判别}
设 $D\subseteq TM$ 是一个光滑分布。则 $D$ 对合当且仅当下列条件成立：
\begin{quote}
若 $\omega$ 是在开子集 $U\subseteq M$ 上零化 $D$ 的任意光滑 $1$-形式，则 $\mathrm{d}\omega$ 也在 $U$ 上零化 $D$。
\end{quote}
\end{proposition}

\begin{proof}
首先，设 $D$ 对合，$\omega$ 是在 $U\subseteq M$ 上零化 $D$ 的光滑 $1$-形式。于是对 $D$ 的任意光滑局部截面 $X,Y$，$\mathrm{d}$ 的公式\eqref{align:ch14-d1form-invariant} 给出
\begin{align*}
\mathrm{d}\omega(X,Y)=X\bigl(\omega(Y)\bigr)-Y\bigl(\omega(X)\bigr)-\omega\bigl([X,Y]\bigr).
\end{align*}
假设蕴含右端每一项在 $U$ 上都为零。

反之，设 $D$ 满足上述条件，$X,Y$ 是 $D$ 的光滑局部截面。若 $\omega^1,\dots,\omega^{n-k}$ 是 $D$ 的光滑局部定义形式，则\eqref{align:ch14-d1form-invariant} 表明对每个 $i=1,\dots,n-k$，
\begin{align*}
\omega^i\bigl([X,Y]\bigr)=X\bigl(\omega^i(Y)\bigr)-Y\bigl(\omega^i(X)\bigr)-\mathrm{d}\omega^i(X,Y)=0,
\end{align*}
这蕴含 $[X,Y]$ 取值于 $D$ 中。故 $D$ 对合。
\end{proof}

正如对合性的 Lie 括号条件只需对一组特定的光滑定义形式加以检验一样，外导数条件也只需在每一点的一个邻域内检验，如下一个命题所示。

\begin{proposition}[对合性的局部余标架判别法]\label{proposition:ch19-局部余标架对合判别}
设 $D$ 是光滑 $n$ 维流形 $M$ 上一个秩为 $k$ 的光滑分布，$\omega^1,\dots,\omega^{n-k}$ 是 $D$ 在开子集 $U\subseteq M$ 上的光滑定义形式。下列各条等价：
\begin{enumerate}[(a)]
\item $D$ 在 $U$ 上对合。
\item $\mathrm{d}\omega^1,\dots,\mathrm{d}\omega^{n-k}$ 都零化 $D$。
\item 存在光滑 $1$-形式 $\bigl(\alpha^i_j\colon i,j=1,\dots,n-k\bigr)$，使得对每个 $i=1,\dots,n-k$，
\begin{align*}
\mathrm{d}\omega^i=\sum_{j=1}^{n-k}\omega^j\wedge\alpha^i_j.
\end{align*}
\end{enumerate}
\end{proposition}

\begin{exercise}\label{exercise:ch19-09}
证明上述命题。
\end{exercise}

再多用一点代数术语，就可以用微分形式给出对合性条件的一种优美表述。回忆我们已把光滑 $n$ 维流形 $M$ 上的光滑微分形式的分次代数定义为 $\Omega^*(M)=\Omega^0(M)\oplus\cdots\oplus\Omega^n(M)$（见第 360 页）。$\Omega^*(M)$ 中的一个\emph{理想}是指一个线性子空间 $I\subseteq\Omega^*(M)$，它对与 $\Omega^*(M)$ 中任意元素的外积封闭；即 $\omega\in I$ 蕴含对每个 $\eta\in\Omega^*(M)$ 都有 $\eta\wedge\omega\in I$。

现在设 $D$ 是光滑 $n$ 维流形 $M$ 上的一个光滑分布。令 $I^p(D)\subseteq\Omega^p(M)$ 表示零化 $D$ 的光滑 $p$-形式空间，并令 $I(D)=I^0(D)\oplus\cdots\oplus I^n(D)\subseteq\Omega^*(M)$。

\begin{exercise}\label{exercise:ch19-10}
对任意光滑分布 $D\subseteq TM$，证明 $I(D)$ 是 $\Omega^*(M)$ 中的一个理想。
\end{exercise}

对应于某个光滑分布 $D$ 的形如 $I(D)$ 的理想有时称为一个\emph{Pfaff 方程组}。若 $\mathrm{d}(I)\subseteq I$，即 $\omega\in I$ 蕴含 $\mathrm{d}\omega\in I$，则称理想 $I\subseteq\Omega^*(M)$ 为一个\emph{微分理想}。

\begin{proposition}[对合性的微分理想判别法]\label{proposition:ch19-微分理想判别}
设 $M$ 是一个光滑流形。一个光滑分布 $D\subseteq TM$ 对合当且仅当 $I(D)$ 是 $\Omega^*(M)$ 中的一个微分理想。
\end{proposition}

\begin{proof}
习题 19-1。
\end{proof}
\end{document}
